% Pursuing Stacks (cotes 134-2 to 134-8): archive cote 134-4, site page 97 (PDF page 98), typescript page 184.
% typescript. Draft compared: ps3.tex:4950-5007.
% AI-transcribed by Claude Opus 5.5 (claude-opus-5-5, adaptive thinking, effort high), Message Batches API, batch ps2 (msgbatch_01NVwr1Yuo8qJf5rykeGBXPK), 2026-11-12; not yet checked by eye.
\scanpage{134-4/97}
\mspage{184}

is the property that F\hand{$\to$} e be a \underline{W}-equivalence, which also means that
the map (**) in (Cat\hand{)} is a \underline{W}-equivalence i.e. is in \underline{W} . If A is \underline{W}-asphe-%
ric, this \unsure{t}\hand{h}ird notion however reduces to the first one, namely F to
be \underline{W}-aspheric.

\uline{T\hand{h}ese terminological conventions,} \uline{in the all-important cases of}
(Cat) \uline{and categories of the type A\textasciicircum,} \uline{should be viewed as the basic ones}\hand{,}
and there should be no question \deleted{here}\inserted{whatever} to requestion them, because\inserted{say} of the
need we \xout{max} are in to dev\over{a}{i}se a terminology, applicable to the general
case of a category M \unsure{endowed} with any saturated set of maps $W \subset F\ell(M)$\hand{,}
\hand{(}which are being thought of as still more general substitutes of ``weak
equivalences''\hand{)}. This shows at once that we will have to renounce to the
name of ``W-aspheric'' which we have used so far, in order to designate
maps which are ``universally in W'' ; indeed, this contradicts the use we
are making of this word, in the case of (Cat). The whole trouble came
from this unappropriate terminology, which slipped in while thinking
of the A\textasciicircum\hand{/}analogy, and forgetting about the still more basic (Cat) !
The mistake is a coarse one indeed, and quite easy to correct - \uline{we}
\uline{better refrain altogether from using the word} ``\uline{aspheric}'' \uline{in the gontext}
\uline{of a general pair} (M,W), \xout{but} and rather speak of maps which are
\unsure{``}\uline{universally in W}'' or ``\uline{universal} W-\uline{equivalences}'', which is indeed more
suggestive, and does not carry any ambiguity. The notion of ``W-aspheric
map'' should be reserved to the case when, among all possible base change
maps Y' \hand{$\to$}Y in M, we can sort out some which we may view as ``localizing
maps'' - a\hand{ll} maps \inserted{I'd think} in case\xout{s} M is a topos, and pretty few ones when
M = (Cat). As for qualifying objects of M, we'll just be specific in
stating properties of the projection X\hand{$\to$}$e_{M}$ - such as being a W-equiva-%
lence, or a universal W-equivalence, or a homotopism for $h_{W}$ (in which
case the name ``contractible object'' is adequate indeed). It may be
convenient, when we got a W, to denote by \xout{UW}
\[ UW \subset W \subset F\ell(M) \]
the corresponding set of maps which are universally in W , \inserted{a property,} which then
can be abbreviated into the simple notation
\[ f \in UW \]
or ``f is in UW''. It should be notes that UW contains all invertible
maps and is stable by composition, but \uline{it} \deleted{is'nt}\inserted{need not be} \uline{saturated}, thus f and \xout{\unsure{x}}
fg may be in UW without g being so.

\margin{97}
